\documentclass[Orbiter User Manual.tex]{subfiles}
\begin{document}

\section{Credits}
%TODO review content
%TODO check websites online
%TODO handle dead websites?

\subsection{Special thanks}
Doug, Josh, Gary, Orb and the entire Orbiter Forum team for keeping things running smoothly, and in particular Josh for providing and maintaining the server for the forum and Orbiter downloads.\\
Jarmo for pushing the envelope with the D3D9 client, and helping debug Orbiter and the graphics interface.\\
All beta testers and bug reporters for their help in getting the new version into shape.\\
All Orbiter users for their continued support. Keep playing!

\subsection{Source contributions and components}
\textbf{XRSound sound module}\\
Copyright (c) Doug Beachy \url{https://www.alteaaerospace.com/}\\
See Sound/XRSound/LICENSE for XRSound license.\\
\\
\textbf{TransX MFD mode}\\
Duncan Sharpe: TransX MFD mode module\\
Steve Arch: TransX development \url{http://orbiter.quorg.org}

\subsection{Libraries, code, data, algorithms}
\textbf{Qt 6}\\
Application framework for the Launchpad, dialogs and windows of the Linux version\\
The Qt Company \url{https://www.qt.io}\\
\\
\textbf{Vulkan}\\
Graphics API of the VulkanClient\\
The Khronos Group \url{https://www.vulkan.org}\\
\\
\textbf{PipeWire}\\
Audio output of XRSound in the Linux version\\
\url{https://pipewire.org}\\
\\
\textbf{dr\_mp3, dr\_flac}\\
MP3 and FLAC decoders for XRSound in the Linux version\\
David Reid \url{https://github.com/mackron/dr_libs}\\
\\
\textbf{stb\_vorbis}\\
Ogg Vorbis decoder for XRSound in the Linux version\\
Sean Barrett \url{https://github.com/nothings/stb}\\
\\
\textbf{libxmp-lite}\\
Tracker module (MOD, S3M, XM, IT) decoder for XRSound in the Linux version\\
\url{https://github.com/libxmp/libxmp}\\
\\
\textbf{libpng}\\
PNG image library\\
\url{http://www.libpng.org/pub/png/libpng.html}\\
\\
\textbf{Zlib}\\
Compression/decompression library\\
Copyright (C) 1995-2013 Jean-loup Gailly and Mark Adler\\
Jean-loup Gailly \href{mailto:jloup@gzip.org}{jloup@gzip.org}\\
Mark Adler \href{mailto:madler@alumni.caltech.edu}{madler@alumni.caltech.edu}\\
\\
\textbf{Font Awesome}\\
Icons font\\
Fonticons, Inc. \url{https://fontawesome.com}\\
CC BY 4.0 License \url{https://creativecommons.org/licenses/by/4.0/}\\
\\
\textbf{NVIDIA Blast 5.1.0}\\
Structural fracture of the Collision module (the low-level and stress solver parts of the SDK)\\
Copyright (c) 2016-2024 NVIDIA Corporation, BSD 3-Clause License (Doc/Licenses/NVIDIA-Blast.txt)\\
\url{https://github.com/NVIDIAGameWorks/Blast}\\
\\
\textbf{VSOP87}\\
Planetary perturbation terms for Mercury to Neptune\\
Bureau des Longitudes, CNRS URA 707\\
P. Bretagnon \href{mailto:pierre@bdl.fr}{pierre@bdl.fr}\\
G. Francou \href{mailto:francou@bdl.fr}{francou@bdl.fr}\\
\\
\textbf{Lunar Solution ELP 2000-82B}\\
Semi-analytical lunar ephemerides\\
Bureau des Longitudes, CNRS URA 707\\
75014, Paris, France\\
M. Chapront-Touze, J. Chapront\\
Astron. Astrophys. 124, 50 (1983)\\
Astron. Astrophys. 190, 342 (1988)\\
\\
\textbf{TASS 1.7}\\
Saturn satellite theory for Mimas, Enceladus, Tethys, Dione, Rhea, Titan, Hyperion and Iapetus\\
A. Vienne, L. Duriez, Astron. Astrophys. 297, 588 (1995)\\
L. Duriez, A. Vienne, Astron. Astrophys. 324, 366 (1997) (Hyperion)\\
\\
\textbf{Earth precession parameters}\\
IAU SOFA C Library\\
\url{http://www.iausofa.org/}\\
\\
\textbf{Planetary precession parameters}\\
IAU/IAG Working Group\\
Report of the IAU/IAG Working Group on cartographic coordinates and rotational elements 2006, \url{http://www.springerlink.com/content/e637756732j60270/}\\
Spin axes and prime meridians of the Sun, Jupiter, Uranus, Neptune, Vesta and 17 moons fitted in the Linux version to the IAU WGCCRE rotation model as NAIF's pck00011 gives it: Archinal et al., "Report of the IAU Working Group on Cartographic Coordinates and Rotational Elements: 2015", Celest. Mech. Dyn. Astr. 130, 22 (2018), and its Correction, Celest. Mech. Dyn. Astr. 131, 61 (2019); all 31 bodies with an IAU model are within 0.2° (pole) and 0.5° (prime meridian) of pck00011 over 1950-2050\\
\\
\textbf{Phobos, Deimos, Miranda, Ariel, Umbriel, Titania, Oberon and Triton ephemeris modules}\\
64-bit modules of the Linux version: precessing Kepler orbits, with periodic terms for Miranda, Ariel, Titania and Oberon\\
Fitted to JPL Horizons over 1800-2200: Mars satellites mar099, Uranus satellites ura184\_merged, Neptune satellites nep098\_merged\\
JPL Solar System Dynamics \url{https://ssd.jpl.nasa.gov/horizons/}\\
Largest position difference from Horizons at 15,000 random epochs not used in the fit (10,000 over 1800-2200 and 5,000 more over 1990-2040), 2000-2040 / 1800-2200: Phobos 7 / 31 km, Deimos 128 / 209 km, Miranda 227 / 1,048 km, Ariel 275 / 1,008 km, Umbriel 870 / 1,270 km, Titania 1,069 / 1,624 km, Oberon 1,039 / 1,444 km, Triton 93 / 497 km\\
\\
\textbf{Proteus, Nereid and Vesta ephemeris modules}\\
Proteus: a precessing Kepler orbit with periodic terms, fitted to JPL Horizons (nep098\_merged) over 1800-2200, within 20 km\\
Nereid and Vesta: Chebyshev tables of JPL Horizons ephemerides (Nereid nep098\_merged, 1800-2199; Vesta JPL\#36, 1600-2500), within 2.5 and 3 km of the Horizons data used; fitted orbits outside the tables\\
\\
\textbf{Pluto system ephemeris modules}\\
Pluto, Charon, Styx, Nix, Kerberos and Hydra: series fitted to JPL's DE441 and plu060 over 1800-2200, within 100 m of them at epochs not used in the fit; two-body motion outside those dates\\
Spin of Pluto and Charon from pck00011; poles and periods of the four small moons: Weaver et al., "The small satellites of Pluto as observed by New Horizons", Science 351, aae0030 (2016)\\
Shapes of Nix and Hydra: Simon Porter, shape models from the New Horizons images (2021), figshare, CC BY 4.0, doi:10.6084/m9.figshare.12779948 and doi:10.6084/m9.figshare.12779975; Kerberos and Styx shaped from the sizes in Weaver et al. (2016)\\
\\
\textbf{Atmosphere profiles}\\
Jupiter: Galileo probe descent, Seiff et al., J. Geophys. Res. 103, 22857 (1998), and the NASA Jupiter fact sheet\\
Saturn: Cassini CIRS and radio occultations, the UVIS thermosphere of Koskinen et al. (2013), and the NASA Saturn fact sheet\\
Uranus and Neptune: Voyager 2 radio occultation and ultraviolet spectrometer, with the NASA fact sheets\\
Titan: Huygens HASI descent profile, Fulchignoni et al., Nature 438, 785 (2005), and Cassini\\
Io: SO\textsubscript{2} atmosphere review, Lellouch et al., in Io after Galileo, Springer (2007)\\
Triton: Voyager 2 and the stellar occultations of 2017 to 2022 (Oliva et al. 2022)\\
Pluto: New Horizons, Gladstone et al., Science 351, aad8866 (2016), and Hinson et al., Icarus 290, 96 (2017)\\
\\
\textbf{Carl Romanik ("Chode")}\\
Semi-major axes in the Phobos, Deimos and Proteus configuration files, the rotation periods of Nereid and Hyperion, and the positions of Venera 9 to 12 and the Pioneer Venus night probe on Venus\\
His 32-bit Phobos, Deimos, Miranda, Ariel, Umbriel, Titania, Oberon and Triton modules, which upstream Orbiter keeps as Windows binaries, are replaced in the Linux version by the modules above
\\
\subsubsection{Harmonic Gravity Models}
\textbf{Mercury Harmonic Gravity Model}\\
MESS160A	\textit{jgmess\_160a\_sha.tab}\\
160x160 spherical harmonic gravity model produced from Magellan radio tracking data.\\
NASA JPL\\
PDS Geosciences Node, Washington University in St. Louis\\
\url{https://pds-geosciences.wustl.edu/dataserv/gravity_models.htm}\\
\\
\textbf{Venus Harmonic Gravity Model}\\
\textit{mod\_shgj120p.a01}\\
120x120 Spherical harmonic gravity model produced from Magellan radio tracking data, modified with initial C(1,0), S(1,0), C(1,1), S(1,1), coefficients padded with zeros for used with Orbiter parser.\\
NASA JPL\\
PDS Geosciences Node, Washington University in St. Louis\\
\url{https://pds-geosciences.wustl.edu/dataserv/gravity_models.htm}\\
\\
\textbf{Earth Harmonic Gravity Model}\\
\textit{egm96\_to360.tab}\\
Modified by Ajaja for use with Orbiter parser.\\
The Development of the Joint NASA GSFC and the National Imagery and Mapping Agency (NIMA) Geopotential Model EGM96, F.G. Lemoine et al, Goddard Space Flight Center, Greenbelt, Maryland, July 1998.\\
\\
\textbf{Moon Harmonic Gravity Model}\\
165x165 Spherical harmonic gravity model produced from radio tracking of the Lunar Orbiters 1‒5, Apollo 15 and 16 subsatellites, Clementine, and the Lunar Prospector spacecraft.\\
NASA JPL\\
PDS Geosciences Node, Washington University in St. Louis\\
\url{https://pds-geosciences.wustl.edu/dataserv/gravity_models.htm}\\
\\
\textbf{Mars Harmonic Gravity Model}\\
MRO120F	\textit{jgmro\_120f\_sha.tab}\\
120x120 Spherical harmonic gravity model, produced from Mars Reconnaissance Orbiter radio tracking data.\\
NASA JPL\\
PDS Geosciences Node, Washington University in St. Louis\\
\url{https://pds-geosciences.wustl.edu/dataserv/gravity_models.htm}\\
\\
\textbf{Vesta Harmonic Gravity Model}\\
"Konopliv, A.S., Park, R.S., Asmar, S.W., and Buccino, D.R., Dawn Vesta Derived Gravity Data, NASA Planetary Data System, DAWN-A-RSS-5-VEGR-V2.0, 2017."\\
NASA JPL\\
\url{https://sbn.psi.edu/pds/resource/dawn/dwnvgravL2.html?refUrl=https\%3A\%2F\%2Fsbn.psi.edu\%2Fpds\%2Farchive\%2Fgravity.html&refName=Gravity&type=Data+Type&typeUrl=https\%3A\%2F\%2Fsbn.psi.edu\%2Fpds\%2Farchive\%2Fdata-types.html}\\
\\


\subsection{Data sources planetary textures}
\textbf{Mercury}\\
Mosaics created using MESSENGER orbital images released by NASA's Planetary Data System (PDS) (\url{http://pds-geosciences.wustl.edu/missions/messenger/index.htm}) on September 7, 2012.\\
NASA/Johns Hopkins University Applied Physics Laboratory/Carnegie Institution of Washington\\
\url{http://messenger.jhuapl.edu/the_mission/mosaics.html}\\
\\
\textbf{Mercury surface labels}\\
USGS\\
Astrogeology Research Program\\
Planetary Geomatics Group\\
Gazetteer of Planetary Nomenclature\\
\url{http://planetarynames.wr.usgs.gov/}\\
\\
\textbf{Venus surface}\\
Composite of Magellan synthetic aperture radar mosaics.\\
Jet Propulsion Laboratory Multimission Image Processing Laboratory\\
Solar System Visualization Project and Magellan science team\\
p45187.tif (5120x2560)\\
\\
\textbf{Venus clouds}\\
Björn Jónsson\\
\url{http://www.mmedia.is/~bjj}\\
\\
\textbf{Earth land surface}\\
Custom processed from Landsat 7 ETM orthorectified imagery.\\
\\
\textbf{Florida}\\
Digital Orthoimage Quarter Quads,\\
Florida Department of Environmental Protection\\
Land Boundary Information System (\url{www.labins.org})\\
Download: \url{ftp://146.201.97.137/DOQQ/2004/RGB/UTM/MrSid}\\
\\
\textbf{Earth water surface}\\
Based on NASA Visible Earth Blue Marble maps\\
\url{http://visibleearth.nasa.gov/}\\
\\
\textbf{Earth night lights}\\
Custom processed, based on NASA Visible Earth Blue Marble maps\\
\url{http://visibleearth.nasa.gov/}\\
\\
\textbf{Earth clouds}\\
NASA Visible Earth Blue Marble maps\\
\url{http://visibleearth.nasa.gov/}\\
\\
\textbf{Earth elevation}\\
SRTM 90m Digital Elevation Data by NASA, released by USGS\\
CGIAR-CSI version 4 Processed for void removal by International Centre for Tropical Agriculture (CIAT)\\
\\
\textbf{Moon surface}\\
LRO LROC-WAC Global Mosaic 100m June2013\\
Arizona State University\\
Astrogeology Science Center\\
\\
\textbf{Moon elevation}\\
LOLA-GDR/Cylindrical\\
\url{http://imbrium.mit.edu/DATA/LOLA_GDR/CYLINDRICAL/IMG/}\\
\\
\textbf{Mars surface}\\
Custom processed, based on NASA MGS/MOC. 256 ppd/230m Mars Odyssey THEMIS-IR Day Global Mosaic 100m v12\\
\url{http://astrogeology.usgs.gov/search/map/Mars/Odyssey/THEMIS-IR-Mosaic-ASU/Mars_MO_THEMIS-IR-Day_mosaic_global_100m_v12}\\
Viking MDIM2.1 Colorized Global Mosaic 232m\\
\url{http://astrogeology.usgs.gov/search/details/Mars/Viking/MDIM21/Mars_Viking_MDIM21_ClrMosaic_global_232m/cub}\\
\\
\textbf{Mars elevation}\\
MOLA Mars elevation data at 128 pixels per degree\\
\url{http://pds-geosciences.wustl.edu/mgs/mgs-m-mola-5-megdr-l3-v1/mgsl_300x/meg128/}\\
\\
\textbf{Mars surface labels}\\
USGS\\
Astrogeology Research Program\\
Planetary Geomatics Group\\
Gazetteer of Planetary Nomenclature\\
\url{http://planetarynames.wr.usgs.gov/}\\
\\
\textbf{Vesta surface}\\
Dawn FC HAMO Global Mosaic 60mp\\
\url{http://astrogeology.usgs.gov/search/details/Vesta/Dawn/DLR/HAMO/Vesta_Dawn_FC_HAMO_Mosaic_Global_74ppd/cub}\\
\\
\textbf{Vesta elevation}\\
Dawn HAMO DTM Global 93mp\\
\url{http://astrogeology.usgs.gov/search/details/Vesta/Dawn/DLR/HAMO/Vesta_Dawn_HAMO_DTM_DLR_Global_48ppd/cub}\\
\\
\textbf{Pluto and Charon surface}\\
New Horizons LORRI/MVIC global mosaics, 300 m/px (NASA/Johns Hopkins University Applied Physics Laboratory/Southwest Research Institute/Lunar and Planetary Institute); Schenk et al., Icarus 314 and 315 (2018); colour from the New Horizons MVIC global colour mosaics of Pluto and Charon (PDS Small Bodies Node, New Horizons Pluto system composition collection, doi:10.26007/mc7j-ef52) on the 300 m mosaics\\
\url{https://planetarymaps.usgs.gov/mosaic/}\\
\\
\textbf{Nix and Hydra surface}\\
Askaniy Anpilogov, Nix and Hydra true color maps (2022) from New Horizons LORRI and MVIC images on the shape models of S. B. Porter et al., Wikimedia Commons, CC BY-SA 3.0; the textures made from them (Nix\_map.dds, Hydra\_map.dds) are under the same licence. Kerberos and Styx have no maps\\
\url{https://commons.wikimedia.org/wiki/File:Nix_true_color_map.png}\\
\url{https://commons.wikimedia.org/wiki/File:Hydra_true_color_map.png}\\
\\
\textbf{Jupiter}\\
"Cassini's best map of Jupiter"\\
NASA/JPL/Space Science Institute\\
Cassini Imaging Central Laboratory for Operations\\
\url{http://www.ciclops.org/view/1270/Cassinis-Best-Maps-of-Jupiter}\\
Rolf Keibel: Jupiter cloud map based on CICLOPS maps\\
\\
\textbf{Io Surface}\\
Based on:\\
Io Galileo SSI/Voyager Color Merged Global Mosaic 1km\\
Astrogeology Science Center\\
USGS\\
\url{http://astrogeology.usgs.gov/search/map/Io/Voyager-Galileo/Io_GalileoSSI-Voyager_Global_Mosaic_ClrMerge_1km}\\
\\
\textbf{Io Elevation}\\
Based on:\\
Oliver L. White, Paul M. Schenk, Francis Nimmo, Trudi Hoogenboom, "A new stereo topographic map of Io: Implications for geology from global to local scales", Journal of Geophysical Research: Planets 119(6), 1276-1301 (2014), doi: 10.1002/2013JE004591\\
\url{http://onlinelibrary.wiley.com/doi/10.1002/2013JE004591/abstract}\\
\\
\textbf{Europa}\\
Based on:\\
Europa Voyager and Galileo SSI Global Mosaic 500m\\
Astrogeology Science Center\\
USGS\\
\url{http://astrogeology.usgs.gov/search/map/Europa/Voyager-Galileo/Europa_Voyager_GalileoSSI_global_mosaic_500m}\\
\\
\textbf{Ganymede}\\
Ganymede Voyager and Galileo Color Global Mosaic 1.4km\\
Astrogeology Science Center\\
USGS\\
\url{http://astrogeology.usgs.gov/search/map/Ganymede/Voyager-Galileo/Ganymede_Voyager_GalileoSSI_Global_ClrMosaic_1435m}\\
\\
\textbf{Callisto}\\
Based on:\\
Callisto Galileo/Voyager Global Mosaic 1km\\
Astrogeology Science Center\\
USGS\\
\url{http://astrogeology.usgs.gov/search/map/Callisto/Voyager-Galileo/Callisto_Voyager_GalileoSSI_global_mosaic_1km}\\
\\
\textbf{Planet textures in the Linux texture download}\\
Ganymede: Askaniy Anpilogov, USGS, Björn Jónsson and NASA/JPL-Caltech/SwRI/MSSS/Brian Swift, Ganymede map (2019); Callisto: Askaniy Anpilogov and John van Vliet, Callisto map (2019); Saturn: Askaniy Anpilogov and Björn Jónsson, Saturn map; Wikimedia Commons, CC BY-SA 3.0, and the tiles made from them are under the same licence\\
Io and Europa from the USGS mosaics above; Europa's colour follows its brightness, fitted on the Galileo colour image PIA19048 (NASA/JPL-Caltech/SETI Institute)\\
Mercury: LobedHomunculus, Mercury texture map from the USGS MESSENGER MDIS global colour mosaic, Wikimedia Commons, CC BY 4.0\\
Uranus: Solar System Scope, CC BY 4.0\\
Triton: Paul Schenk, Lunar and Planetary Institute, Voyager 2 global map (NASA/JPL-Caltech/LPI), on the earlier Orbiter map where Voyager saw nothing\\
\url{https://commons.wikimedia.org/wiki/File:Ganymede_map_by_Askaniy.png}\\
\url{https://commons.wikimedia.org/wiki/File:Callisto_map_by_Askaniy.png}\\
\url{https://commons.wikimedia.org/wiki/File:Saturn_map_by_Askaniy.png}\\
\url{https://photojournal.jpl.nasa.gov/catalog/PIA19048}\\
\url{https://commons.wikimedia.org/wiki/File:Mercury_Texture_Map.png}\\
\url{https://commons.wikimedia.org/wiki/File:Solarsystemscope_texture_2k_uranus.jpg}\\
\url{https://commons.wikimedia.org/wiki/File:Triton_map_no_grid.jpg}\\
\\
\textbf{Io, Europa, Ganymede, Callisto surface labels}\\
USGS\\
Astrogeology Research Program\\
Planetary Geomatics Group\\
Gazetteer of Planetary Nomenclature\\
\url{http://planetarynames.wr.usgs.gov/}\\
\\
\textbf{Saturn}\\
Björn Jónsson\\
\url{http://www.mmedia.is/~bjj}\\
Rolf Keibel: texture adaptation\\
\\
\textbf{Saturn rings}\\
Björn Jónsson\\
\url{http://www.mmedia.is/~bjj}\\
\\
\textbf{Mimas}\\
NASA/JPL-Caltech/SSI/Lunar and Planetary Institute\\
Cassini Imaging Central Laboratory for Operations\\
PIA 18437\\
\url{http://www.ciclops.org/view/7963/Color-Maps-of-Mimas---November-2014}\\
\\
\textbf{Enceladus}\\
NASA/JPL-Caltech/SSI/Lunar and Planetary Institute\\
Cassini Imaging Central Laboratory for Operations\\
PIA 18435\\
\url{http://www.ciclops.org/view/7961/Color-Maps-of-Enceladus---November-2014}\\
\\
\textbf{Tethys}\\
NASA/JPL-Caltech/SSI/Lunar and Planetary Institute\\
Cassini Imaging Central Laboratory for Operations\\
PIA 18439\\
\url{http://www.ciclops.org/view/7965/Color-Maps-of-Tethys---November-2014}\\
\\
\textbf{Dione}\\
NASA/JPL-Caltech/SSI/Lunar and Planetary Institute\\
Cassini Imaging Central Laboratory for Operations\\
PIA 18434\\
\url{http://www.ciclops.org/view/7960/Color-maps-of-Dione---November-2014}\\
\\
\textbf{Rhea}\\
NASA/JPL-Caltech/Space Science Institute/Lunar and Planetary Institute\\
\url{http://photojournal.jpl.nasa.gov/catalog/PIA18438}\\
\\
\textbf{Titan surface}\\
NASA/JPL/Space Science Institute/Cassini Data Analysis Program/USGS Astrogeology Science Center/Ian Regan\\
Updated, amended and restored version of USGS photomosaic\\
\url{https://astrogeology.usgs.gov/search/map/Titan/Cassini/Global-Mosaic/Titan_ISS_P19658_Mosaic_Global_4km}\\
\url{https://astrogeology.usgs.gov/search/map/Titan/Cassini/Global-Mosaic/Titan_ISS_Globe_65Sto45N_450M_AvgMos}\\
\url{https://www.flickr.com/photos/10795027@N08/43023455582/}\\
\url{https://www.insaturnsrings.com/titan-seam-blending}\\
Map used with permission\\
\\
\textbf{Titan elevation}\\
P. Corlies, A. G. Hayes, S. P. D. Birch, R. Lorenz, B. W. Stiles, R. Kirk, V. Poggiali, H. Zebker, L. Iess "Titan's Topography and Shape at the End of the Cassini Mission", Geophysical Research Letters 44(23), 11754-11761 (2017)\\
\url{https://agupubs.onlinelibrary.wiley.com/doi/full/10.1002/2017GL075518}\\
\\
\textbf{Titan surface labels}\\
USGS\\
Astrogeology Research Program\\
Planetary Geomatics Group\\
Gazetteer of Planetary Nomenclature\\
\url{http://planetarynames.wr.usgs.gov/}\\
\\
\textbf{Iapetus}\\
NASA/JPL-Caltech/SSI/Lunar and Planetary Institute\\
Cassini Imaging Central Laboratory for Operations\\
PIA 18436\\
\url{http://www.ciclops.org/view/7962/Color-Maps-of-Iapetus---November-2014}\\
\\
\textbf{Phoebe}\\
NASA/JPL/Space Science Institute\\
Cassini Imaging Central Laboratory for Operations\\
PIA 07775\\
\url{http://www.ciclops.org/view/1743/Map-of-Phoebe---December-2005}\\
\\
\textbf{Uranus}\\
James Hastings-Trew\\
\url{http://apollo.spaceports.com/~jhasting/}\\
Rolf Keibel: texture adaptation\\
\\
\textbf{Miranda, Ariel, Umbriel, Titania, Oberon}\\
Robert Stettner\\
Credits: Planetary Satellite Mean Orbital Parameters and Moon Maps\\
\\
\textbf{Neptune}\\
James Hastings-Trew\\
\url{http://apollo.spaceports.com/~jhasting/}\\
\\
\textbf{Triton, Proteus, Nereid}\\
Robert Stettner\\
Credits: Planetary Satellite Mean Orbital Parameters and Moon Maps\\
Rolf Keibel: Triton texture adaptation from Voyager images


\subsection{Celestial sphere}
\textbf{Background star databases}\\
ESA, ed. 1997, The Hipparcos and Tycho Catalogues, SP No. 1200 (ESA)\\
\url{https://www.cosmos.esa.int/web/hipparcos/hipparcos-2}\\
\url{http://cdsarc.u-strasbg.fr/ftp/cats/I/239/}\\
F. van Leeuwen (2005). A new reduction of the raw Hipparcos data, Astron. Astrophys. 439(2): 791-803. 10.1051/0004-6361:20053193\\
\url{https://www.researchgate.net/publication/1760034_Validation_of_the_new_Hipparcos_reduction}\\
\\
\textbf{Background map: Starmap 2020}\\
NASA/Goddard Space Flight Center Scientific Visualization Studio.\\
Gaia DR2: ESA/Gaia/DPAC.\\
\url{https://svs.gsfc.nasa.gov/4851}\\
\\
\textbf{Background maps: DDS2 (visible), Hydrogen alpha, IRAS (far IR), Planck (Microwave, Source: ESA/Planck), Radio, RASS (X-ray), Fermi (Gamma)}\\
Chromoscope \url{http://www.chromoscope.net/}\\
Stuart Lowe, Chris North (Cardiff University) and Robert Simpson (Oxford University)\\
\\
%TODO name doesn't match what is shown in launchpad
\textbf{Background map: WMAP Microwave images}\\
WMAP Science Team\\
WMAP "Science on a sphere" microwave sky images\\
NASA/LAMBDA\\
\url{http://lambda.gsfc.nasa.gov/product/map/current/sos/}\\
\\
\textbf{Constellation boundaries}\\
(c) 2012-2026 Pierre Barbier,\\
Interpolated merged edges (J2000)\\
Creative Commons Attribution-ShareAlike 4.0 International License\\
\url{https://pbarbier.com/constellations/boundaries.html}\\
\url{https://pbarbier.com/constellations/lines_in_20.txt}\\


\subsection{Spacecraft and structure models and textures}
\textbf{DeltaGlider and DG-S mesh and virtual cockpit}\\
Roger "Frying Tiger" Long\\
\\
\textbf{Space Shuttle Atlantis}\\
Michael Grosberg: meshes and textures\\
Don Gallagher: mesh and texture extensions\\
Robert Conley ("estar"): Module extensions: Movable arm and grappling, including MMU and Satellite extensions; documentation\\
David Hopkins: Module code extensions\\
Damir Gulesich: Space Shuttle External Tank and Solid Rocket Booster mesh and textures.\\
\\
\textbf{LDEF mesh and textures}\\
Don Gallagher\\
\\
\textbf{ISS model "Project Alpha"}\\
Andrew Farnaby\\
\\
\textbf{Mir model}\\
Jason Benson ("agent036")\\
\\
\textbf{Dragonfly model}\\
Roger "Frying Tiger" Long: Mesh improvements and textures\\
Radu Poenaru: Electrical and environmental simulation, Dragonfly panels\\
\\
\textbf{Shuttle-A model}\\
Roger "Frying Tiger" Long: Shuttle-A mesh\\
Radu Poenaru: Virtual cockpit and cargo management\\
\\
\textbf{Hubble Space Telescope (HST) model}\\
David Sundstrom\\
\\
\textbf{KSC VAB mesh}\\
Valerio Oss\\
\\
\textbf{PTV (Personal transport vehicle) mesh}\\
Balázs Patyi\\
\href{mailto:patyibalazs@yahoo.com}{patyibalazs@yahoo.com}\\
\\
\textbf{Default exhaust texture, cloud microtextures}\\
"McWgogs"\\
\url{http://mcwgogs.deviantart.com/}


\end{document}

